\section{Results}

\subsection{Correctness and Availability}

Table~\ref{tab:correctness} reports the three correctness workloads. In the forced schedule, all 10 follow-up GETs in C returned the obsolete value after the delayed refill. A and D returned the committed availability in all 10 follow-up GETs per mode. D counted one rejected refill per trial. The original delayed GET still returned the older SQL snapshot in every mode. The distinction matters because the correctness property concerns what the delayed reader leaves in Redis for a subsequent customer.

\begin{table}[t]
\centering
\caption{Correctness outcomes across ten trials per reported mode and workload. Counts are pooled for description. Browsers within each stockout round share a single sale.}
\label{tab:correctness}
\setlength{\tabcolsep}{4pt}
\begin{tabular}{@{}c r r r@{}}
\toprule
Mode & \shortstack{Stale forced\\follow-up} & \shortstack{False-positive\\checkout} & \shortstack{Stale hot-item\\GETs} \\
\midrule
A & 0/10 & 0/80 & 0/34,755 \\
B & --- & 80/80 & 36,069/36,169 \\
C & 10/10 & 0/80 & 1,339/35,976 \\
D & 0/10 & 0/80 & 0/35,498 \\
\bottomrule

\end{tabular}
\end{table}

In the stockout trial dataset, eight browsers per round could attempt to order only if the one-portion display remained. TTL-only B produced 80 attributable stockout rejections in 80 browsing sessions. A, C, and D produced none in their 80 sessions each. Their browsers saw zero available portions and never issued the checkout. The 80 decisions in each mode are grouped in ten shared-inventory trials. Each trial passed the inventory audit, and only the two initiating sales committed. The test confirms that a correct database rejection can still follow a misleading item display.

In the ungated single-item stress test, every order updated the ingredient revisions of the item all six readers were requesting. Plain deletion returned 1{,}339 stale responses across 35{,}976 eligible GETs. Its median stale share across the ten trials was 2.56\%, with an observed range of 0.88\% to 12.31\%. D recorded 339 rejected generation checks and zero stale responses across 35{,}498 eligible GETs. TTL-only B had the highest stale count, consistent with a three-second TTL in a 1.5-second workload with frequent writes. The counts in Table~\ref{tab:correctness} describe eligible successful GETs and summarize ten trials per mode.


\subsection{Arrival-Rate Screen and Repeated Comparison}

Fig.~\ref{fig:screen} shows the screening pass. All four modes met the screening criterion up to 2{,}400 requests/s. In its single 4{,}800 requests/s screen, D delivered 87.36\% of offered requests and recorded 1{,}810 dropped iterations. This selected the rate for repeated confirmation. At 8{,}000 requests/s, all modes missed the 95\% delivery threshold. The individual screening runs locate pressure for this host and driver.

\begin{figure}[t]
\centering
\includegraphics[width=\columnwidth]{open_screen}
\caption{Open-loop rate screen. Each point represents one three-second trial. The dotted line marks target delivery.}
\label{fig:screen}
\end{figure}

\begin{table*}[t]
\centering
\caption{Open-loop comparison at a three-second TTL. Each row summarizes ten trials. GET and POST columns are medians of per-trial P95 among successful responses. The criterion requires at least 95\% in-window delivery and at most 50\,ms overall P95.}
\label{tab:performance}
\setlength{\tabcolsep}{6pt}
\begin{tabular}{@{}l c r r r r r c@{}}
\toprule
Offered req/s & Mode & GET P95 (ms) & POST P95 (ms) & Successful req/s & Dropped (\%) & Menu SELECTs/1,000 GETs & Passes \\
\midrule
600 & A & 1.17 & 3.22 & 599.8 & 0.00 & 1000.0 & 10/10 \\
600 & B & 0.99 & 3.67 & 600.0 & 0.00 & 57.9 & 10/10 \\
600 & C & 1.45 & 3.93 & 599.8 & 0.00 & 361.6 & 10/10 \\
600 & D & 1.30 & 3.61 & 599.8 & 0.00 & 362.4 & 10/10 \\
\midrule
4800 & A & 28.67 & 35.31 & 4646.2 & 3.16 & 1000.0 & 9/10 \\
4800 & B & 4.14 & 54.30 & 4777.3 & 0.43 & 10.3 & 10/10 \\
4800 & C & 31.53 & 45.60 & 4681.7 & 2.42 & 424.4 & 10/10 \\
4800 & D & 34.31 & 41.90 & 4599.1 & 4.09 & 524.7 & 7/10 \\
\bottomrule

\end{tabular}
\end{table*}

\begin{table*}[t]
\centering
\caption{Paired changes at 4{,}800 offered requests/s. Differences are candidate minus baseline within each repeat, summarized by the median and its 95\% percentile bootstrap interval.}
\label{tab:paired}
\begin{tabular}{@{}lrrr@{}}
\toprule
Contrast & GET P95 (ms) & POST P95 (ms) & Delivery (percentage points)\\
\midrule
B minus A & -23.96 (-26.28, -22.71) & +20.53 (+2.45, +29.44) & +2.37 (+1.74, +3.22) \\
D minus C & +3.25 (-1.41, +15.50) & -2.27 (-5.71, +6.51) & -1.57 (-3.23, -0.99) \\
\bottomrule

\end{tabular}
\end{table*}

\begin{figure}[t]
\centering
\includegraphics[width=\columnwidth]{open_latency}
\caption{Median per-trial endpoint P95 at matched offered rates. Whiskers show 95\% percentile bootstrap intervals over trials.}
\label{fig:latency}
\end{figure}

Table~\ref{tab:performance} gives the repeated comparison, and Fig.~\ref{fig:latency} shows the variation in GET and POST P95 separately. The 600 requests/s comparison met the specified criterion in all ten trials per policy. Its median in-window delivery was at least 99.9\% for every policy.

At 4{,}800 offered requests/s, B had the lowest median item GET P95, 4.14\,ms, and the highest median goodput, 4{,}777.3 successful requests/s. Its median POST P95 was 54.30\,ms, compared with 35.31\,ms in A. Table~\ref{tab:paired} shows the paired changes. The POST difference for B relative to A had the opposite sign from the GET difference. The load criterion uses overall P95, so an individual endpoint can exceed 50\,ms even when a trial passes. A passed nine of ten trials, B and C passed ten each, and D passed seven. The recorded load shortfalls came from scheduled starts that k6 dropped and responses that finished outside the measurement window. The repeated open-loop trials recorded no failed or incomplete HTTP attempts after a start.

At 4{,}800 requests/s, D used a median 524.7 item SELECTs per 1{,}000 started GETs, compared with 424.4 in C and 1{,}000 in A. D counted a median of 3 rejected generation checks per trial at 600 requests/s and 3{,}417.5 at 4{,}800 requests/s. A generation mismatch shows a refill crossed an invalidation. It can conservatively reject a candidate that has already seen the newer commit. The paired median change in delivery from C to D was about \(-1.6\) percentage points. The paired GET and POST P95 intervals for that comparison spanned zero. The data support a delivery cost for D within this test setup. They leave the separate effect of Redis command work on latency unresolved. Omitting repeat 1 preserved the ordering of the four delivery medians.

In the TTL sweep (Fig.~\ref{fig:ttl}), B's median payload hit rate increased from 86.7\% at 1\,s to 98.1\% at 10\,s. Its median GET P95 fell from 1.04 to 0.74\,ms. D's median hit rate stayed between 62.4\% and 63.8\%, with median GET P95 near 1.2\,ms at all three TTLs. That weaker TTL sensitivity is consistent with orders removing payloads before expiry. The sweep measured hit rate and response time at 600 requests/s.


\begin{figure}[t]
\centering
\includegraphics[width=\columnwidth]{ttl_sweep}
\caption{TTL sensitivity at 600 offered requests/s. Points show medians from ten 12-second trials per cell. Whiskers show 95\% percentile bootstrap intervals over trials.}
\label{fig:ttl}
\end{figure}
